% Values: PRECISE, FAITHFUL, ILLUMINATING, ECONOMICAL
% Evidence planned:
% - PRECISE: a rule for every bare production of AtExp, ExpRow, Exp, Match,
%   MRule, Dec, ValBind, AtPat, PatRow and Pat (book plan "Bare Constructs",
%   with Ty and TyRow removed by the reduced syntax); every object and
%   operation defined before its first use (Val, Rec, FV, APPLY as a finite
%   set); one place, the group "Application of Values", defining application.
% - FAITHFUL: number specialisation and typed integer constants (A87, A27),
%   from the principal derivation (A95);
%   no state and no exceptions (D2); total arithmetic through APPLY
%   (D3, A40, A9); APPLY as a set for function equality (A51) and Debug.todo
%   (A62), Debug.log returns its argument (A62); premise order carries no
%   meaning (A57, A58); strict prefix (&&) and (||) (A56); closures capture
%   their own call's values (A63); no rule for a forced delayed
%   self-reference (A25) or resource exhaustion (A61).
% - ILLUMINATING: comments on why APPLY is a set, why two rules cover Dec,
%   the absent rules for failed matches, strictness of let, the order of
%   match rules, and why constructor application has no arity side
%   condition.
% - ECONOMICAL: options merge tuples, records and binding sequences; derived
%   forms carry no rules; equality and ordering are left to APPLY.
\chapter{Dynamic Semantics for the Core}\label{ch:dynamic-core}

The dynamic semantics of the Core defines the values that phrases evaluate
to, the environments that declarations and patterns produce, and the closures
by which functions are represented. Its sentences infer that a phrase
evaluates to a value or an environment in a given environment, or that a
pattern, matched against a value, gives an environment or fails. Evaluation
in Standard Elm has no state and no exceptions: sentences carry no states,
and no rule produces an exception packet (Section~\ref{sec:dynamic-core-rules}).
% D2
Every phrase evaluated here is assumed to have elaborated in the static
semantics of Chapter~\ref{ch:static-core}. % S4

\section{Reduced Syntax}\label{sec:dynamic-core-reduced}

Types play no part in evaluation, so for the dynamic semantics the Core
syntax is reduced as follows:
% D1 (definition-style)
\begin{itemize}
\item The type annotation ``$\kw{:}\ \ty$'' is omitted from every value
  binding, which thus has the form
  $\var\ \kw{=}\ \expm\ \opt{\kw{and}\ \valbind}$ or
  $\atpat\ \kw{=}\ \expm$.
\item The Core phrase classes Ty and TyRow are omitted.
\end{itemize}

\section{Number Specialisation}\label{sec:dynamic-core-numbers}

Erasing types would leave the value of an integer constant undetermined when
its type is Float or a type variable of class $\knum$. The dynamic semantics
of every level is therefore given for the elaborated program after
\emph{number specialisation}, a translation of the whole program, which we
define here, where the dynamic semantics begins. % A87
It starts from the principal derivation of the program
(Section~\ref{sec:static-core-principal}), with the following default.

\emph{Defaulting of number.} In the principal derivation, each type variable
of class $\knum$ that is bound neither by a closure
(rules~\ruleref{r:s-dec-group}, \ruleref{r:sm-topdec-value}
and~\ruleref{r:sm-datbind}), nor by the quantifier $\forall U'$ of an annotated
binding (rule~\ruleref{r:s-valbind-var}), nor as a parameter of a type
declaration (rules~\ruleref{r:sm-datbind}, \ruleref{r:sm-typbind}
and~\ruleref{r:sm-typbind-record}), is taken to be $\tFloat$: $\tFloat$ is
substituted for it throughout the derivation, which remains a derivation for
the program, since $\tFloat \in \Adm{\knum}$. Every constant of such a type is
an integer constant, since a floating-point constant has the type $\tFloat$
itself, and every integer constant may have the type $\tFloat$. % A27, A95
A variable of class $\knum$ bound in any of these ways is not affected: one
bound by a closure or an annotation is quantified in the type scheme of a value,
and the copies below assign it; a parameter of a type declaration is bound by a
type function, whose applications substitute for it.
Since closure never generalises an explicit variable in scope
(Section~\ref{sec:static-core-closure}), an explicit variable is bound by the
annotation that introduces it, and so is never defaulted, even where a nested
binding uses it: with \kw{f :\ number -> number} and
\kw{f x = let g y = y + x in g 1}, the \kw{number} of \kw{x} is bound by the
annotation of \kw{f}, not by the closure of \kw{g}. % A27 (2026-11-10)
The default changes no type that the program forces, and so no type a program
elaborates with: it decides only types the principal derivation leaves open.

Then each declaration of values, of a \kw{let} or at the top level of a
module, whose bindings have type schemes that quantify type variables of class $\knum$ is
replaced by the following copies of it:
\begin{itemize}
\item the copy that assigns $\tFloat$ to each of those variables, the
  \emph{default copy}, which binds the variables of the declaration under
  their own names; and
\item one copy for each other assignment of $\tInt$ or $\tFloat$ to those
  variables at which the program uses the declaration, the copies binding
  distinct new variables.
\end{itemize}
Each occurrence of a variable the declaration binds is made to refer to the
copy for the instance at which that occurrence was elaborated. Every binding
of a declaration quantifies all the variables of class $\knum$ of the
declaration, those its own type does not mention included
(Section~\ref{sec:static-core-closure}), so an instance assigns a type to each
of them; a variable that nothing at the occurrence fixes is $\tFloat$, as a
variable left free in the program is by the default above. % A93
A declaration
whose bindings quantify no type variable of class $\knum$ is not copied, and
is its own default copy. Since the
default copy is made whether or not the program uses the declaration at that
instance, every declaration has a copy: in particular every top-level value
of a module, including \kw{main} and every value that no phrase uses, has a
copy and is evaluated when its module is evaluated. % A59, A87
The default copy of the declaration of \kw{main} is the one that
Chapter~\ref{ch:programs} names by $\pgDflt$, at the instance
$\pgDflt\sigma$ of the scheme $\sigma$ of \kw{main}. % A27
Copies are made from the outside in, so that a variable of class $\knum$
that a declaration does not quantify is already fixed by the copy of the
declaration that encloses it, and a variable of class $\knum$ left free in
the program has been taken to be $\tFloat$ by the default above. % A27
With \kw{check h = if h == 0 then h else
Debug.todo "h"} and \kw{x = check (2 \textasciicircum{} -1)} used only as
\kw{x // 2}, the copy of \kw{x} at $\tInt$, where \kw{2 \textasciicircum{} -1}
is $0$ (Appendix~\ref{ch:appendix-dynamic-basis}), is $0$; but the default
copy, at $\tFloat$, where it is $0.5$, has no value, so the program has none.
A variable of class $\knum$ that
the type of \kw{x} does not mention, such as that of \kw{2 \textasciicircum{}
-1} in \kw{if 2 \textasciicircum{} -1 == 0 then 1 else 2}, is not a variable of
the declaration: no closure generalises it, and it is $\tFloat$ in every copy.
% A87, A43, A27
After specialisation every integer constant has the type $\tInt$ or
$\tFloat$. The specialisation is finite: $\Adm{\knum}$ contains only $\tInt$
and $\tFloat$ (Section~\ref{sec:static-core-classes}), so a declaration with
$k$ variables of class $\knum$ has at most $2^k$ copies, even when a binding
is used recursively at an instance other than its own.
% DECIDED (audit X6, X7, decision 4): every declaration keeps its default
% copy, so that main and unused values are evaluated (A59); for a let
% declaration this keeps the Comment on rule r:d-atexp-let true.

Where the type of an occurrence of an integer constant matters below, it is
its type after specialisation.

\section{Simple Objects}\label{sec:dynamic-core-simple}

The dynamic semantics builds its objects from the identifier classes of
Chapter~\ref{ch:syntax-core} and the simple object classes of
Figure~\ref{fig:dynamic-core-simple}.

\begin{figure}[htbp]
\centering
\begin{objtable}
\objrow{\sv}{\SVal}{special values}
\objrow{b}{\BasVal}{basic values}
\objrow{\con}{\cls{Con}}{constructor names}
& & \FAILcls & & failure\\
\end{objtable}
\caption{Simple Semantic Objects}
\label{fig:dynamic-core-simple}
\end{figure}

SVal is the class of \emph{special values}: the integers of the Int range
$\IntRange$, the floating-point numbers (among them the two infinities and
NaN), the characters, and the strings, which are finite sequences of
characters. The range $\IntRange$ and the representation of characters are
implementation parameters, given in Appendix~\ref{ch:appendix-implementation}.
% D8, A40, A52
Each occurrence of a special constant $\scon$ denotes a special value,
written $\valof(\scon)$. An integer constant spells an integer $n$, and its
value depends on the type of the occurrence after number specialisation
(Section~\ref{sec:dynamic-core-numbers}): at $\tInt$,
$\valof(\scon)=\norm(n)$, where $\norm$ maps an integer into $\IntRange$ by
the same normalisation as Int arithmetic
(Section~\ref{sec:appendix-implementation-int}); % A9, A40
at $\tFloat$, $\valof(\scon)$ is the floating-point number nearest to $n$,
rounded as Appendix~\ref{ch:appendix-dynamic-basis} prescribes. % A87
A floating-point constant denotes the real number it spells, rounded to the
nearest floating-point number; a character constant denotes a character, and
a string constant a sequence of characters, as explained in
Chapter~\ref{ch:syntax-core}.

BasVal is the class of \emph{basic values}: the names of the operations of
the initial basis, listed in Appendix~\ref{ch:appendix-dynamic-basis}. Each
basic value $b$ has an \emph{arity} $\arity(b)\geq 1$, the number of
arguments it takes before it acts (Section~\ref{sec:dynamic-core-basic}).

The class Con of value constructors serves also as the class of
\emph{constructor names}. We take Con to contain, besides the constructors
spelt by identifiers, the two list constructor names \kw{[]} and \kw{::},
which no identifier spells.
% DECIDED: list values are constructor values built from two constructor
% names [] and :: that are not identifiers; the registry's Val has no list
% class of its own, and the operator table of Appendix D gives (::) the
% constructor :: (audit X1, X5).
Each constructor $\con$ has an \emph{arity} $\arity(\con)\geq 0$, the number
of arguments it takes: for a constructor of a custom type, the number of
types that follow it in its declaration (Chapter~\ref{ch:syntax-modules});
$\arity(\kw{[]})=0$ and $\arity(\kw{::})=2$. Constructors of different custom
types may be spelt alike and have different arities; $\arity(\con)$ always
means the arity of $\con$ in the custom type of the value at hand, which
elaboration determines. Since only phrases that elaborate are evaluated, a
constructor value is only ever compared with the constructors of its own
custom type, whose names are distinct, so no confusion arises.
% DECIDED: constructor names are not qualified by their module; two
% constructor values meet in a rule only at one custom type.

FAIL is the outcome of matching a value against a pattern that does not fit
it. It is not a value: it serves only within the rules, to record that a
match has not succeeded.

\section{Compound Objects}\label{sec:dynamic-core-compound}

The compound objects of the dynamic semantics are shown in
Figure~\ref{fig:dynamic-core-compound}. The classes Val, Closure, Env, QualEnv
and ValEnv are defined together; they are the smallest classes closed under
the constructions of the figure.

\begin{figure}[htbp]
\centering\small
\begin{objtable}
\objdef{v}{\Val}{\{()\}\cup\SVal\cup\BasApp\cup\ConVal\cup\Record}
& & & \phantom{{}={}}\cup\Tuple\cup\Closure\cup\EffVal\cup\JVal & \\
\objdef{(b, v_1\cdots v_k)}{\BasApp}{\textstyle\bigcup_{k\geq 0}\BasVal\times\Val^k}
\objdef{\con\ v_1\cdots v_k}{\ConVal}{\textstyle\bigcup_{k\geq 0}\cls{Con}\times\Val^k}
\objdef{r}{\Record}{\cls{Lab}\finmap\Val}
\objdef{(v_1,v_2\opt{,v_3})}{\Tuple}{\Val^2\cup\Val^3}
\objdef{(\atpat,\expm,E,\VE)}{\Closure}{\cls{AtPat}\times\cls{Exp}\times\Env\times\ValEnv}
\objdef{(\QE,\VE)\ \text{or}\ E}{\Env}{\QualEnv\times\ValEnv}
\objdef{\QE}{\QualEnv}{\cls{ModId}\finmap\Env}
\objdef{\VE}{\ValEnv}{(\cls{Var}\cup\cls{Con})\finmap\Val}
\end{objtable}
\caption{Compound Semantic Objects}
\label{fig:dynamic-core-compound}
\end{figure}

The value $()$ is the unit value. A member $(b, v_1\cdots v_k)$ of BasApp is
a basic value applied to the arguments it has received so far, and always has
$k<\arity(b)$; we identify $(b)$ with $b$, so that $\BasVal\subseteq\BasApp$.
A member $\con\ v_1\cdots v_k$ of ConVal is a \emph{constructor value}, and
always has $k\leq\arity(\con)$: when $k=\arity(\con)$ it is complete, and
when $k<\arity(\con)$ it is a constructor awaiting further arguments. A
constructor name $\con$ with $k=0$ is itself a constructor value. We write
$v_1\mathbin{\kw{::}}v_2$ for the constructor value $\kw{::}\ v_1\ v_2$; the
lists are thus the constructor values $\kw{[]}$ and $v_1\mathbin{\kw{::}}v_2$
with $v_2$ a list. Records are finite maps from labels to values; the
empty record is $\emptymap$. A tuple has two or three components.
The classes EffVal of effect values (task, command, subscription,
process-identifier and program values) and JVal of JSON values are defined
in Chapter~\ref{ch:effects}; this chapter uses only the fact that their
members are values. % D10, audit X11

Projection, injection and modification are as defined for the static
semantics (Chapter~\ref{ch:static-core}); the convention that a decorated
metavariable ranges over the class of the undecorated one is stated in
Chapter~\ref{ch:introduction}. % O5, audit X14
The
unions in Figure~\ref{fig:dynamic-core-compound} are disjoint unions, and the
classes defined there are disjoint from one another; the injections into Val
are left implicit. An environment $E$ is modified by a
value environment as $E+\VE$, meaning $E+(\emptymap,\VE)$.
The names $E$, $\VE$, Env and the like denote dynamic objects here and
static objects in Chapter~\ref{ch:static-core}; the two semantics are
presented apart, so the reuse is harmless.

A long identifier is looked up in an environment $E=(\QE,\VE)$ as follows:
\[
  E(\var)=\VE(\var),\qquad
  E(\modid\kw{.}\var)=(\QE(\modid))(\var),
\]
and likewise for constructors. The identifiers that the equivalent forms of
derived forms introduce (Appendix~\ref{ch:appendix-derived-forms}) are looked
up in the same way. Because only an import of the module \kw{Basics} provides
the qualifier \kw{Basics} (Section~\ref{sec:syntax-modules-restrictions}),
\kw{Basics.True}, \kw{Basics.False} and \kw{Basics.negate} always denote the
values of the initial dynamic basis
(Appendix~\ref{ch:appendix-dynamic-basis}); the constructor \kw{::} is bound
in the initial environment $E_0$ of that appendix, from which every module is
evaluated. % audit X18, V3; A89
Operators are not looked up in $E$, and are
not in the domain of a value environment. Since no program declares an
operator, every operator is one of the initial basis, and its value is the
same in every environment: the value $\OTdyn(\opv)$ that the dynamic
operator table of Appendix~\ref{ch:appendix-dynamic-basis}
(Figure~\ref{fig:dynamic-basis-operators}) gives it. For \kw{::} this value
is the list constructor \kw{::}. % A16, audit X1, X5

\section{Basic Values}\label{sec:dynamic-core-basic}

The meaning of the basic values is represented by a function $\APPLY$ which,
for a basic value $b$ of arity $n$ and values $v_1,\ldots,v_n$, gives a finite
set of values:
\[
  \APPLY(b, v_1\cdots v_n)\in\Fin\Val .
\]
$\APPLY$ is defined in Appendix~\ref{ch:appendix-dynamic-basis}
(Section~\ref{sec:appendix-dynamic-basis-apply}), which also gives the arity
of every basic value. Equality and the order relations of the initial basis
are basic values, and their meaning on structured values is defined there,
not in this chapter. % A50, A53

$\APPLY$ gives a set, rather than a value, so that the definition is exact
about the two operations of the initial basis that do not determine one
result. For every basic value the set has exactly one member, with two
exceptions. First, when an equality test reaches two functions its result is
an unspecified Bool, so that the set may hold both \kw{True} and \kw{False};
evaluation is then non-deterministic. % A51
Second, for the basic value $b$ bound to \kw{Debug.todo}, $\APPLY(b,v)$ is
empty, so that an application of \kw{Debug.todo} has no result. % A62
Every other basic value gives a result for every argument: arithmetic gives a
defined value even where no mathematically correct one exists, % D3
and \kw{Debug.log} returns its second argument, any output it makes lying
outside the definition. % A62
Thus no evaluation ends in a run-time error: a phrase either has a value or
has none. % D2, D3

\section{Function Closures}\label{sec:dynamic-core-closures}

A \emph{function closure} $(\atpat,\expm,E,\VE)$ is applied to a value $v$ by
matching $v$ against $\atpat$ and evaluating $\expm$ in the resulting
environment. The environment is not $E+\VE$ but $E+\Rec\VE$, where the
domain of the fourth component $\VE$ holds the variables treated recursively,
and
\[
  \Rec : \ValEnv\to\ValEnv
\]
is defined by:
\begin{itemize}
\item $\Dom(\Rec\VE)=\Dom\VE$;
\item if $\VE(\var)\notin\Closure$, then $(\Rec\VE)(\var)=\VE(\var)$;
\item if $\VE(\var)=(\atpat,\expm,E,\VE')$, then\\
  $(\Rec\VE)(\var)=(\atpat,\expm,E,\VE)$.
\end{itemize}
Thus $\Rec\VE$ gives each closure in $\Ran\VE$ the bindings of $\VE$ afresh,
so that a recursive application made during the evaluation of $\expm$ finds
them; the recursion is unwound one step at a time, and every semantic object
stays finite. $\Rec$ is used in two places in the rules: for a recursive
declaration, and for the application of a closure.

A declaration is recursive when one of its bindings refers to a variable it
binds. For a value binding $\valbind$, $\dcoreFV(\valbind)$ is the set of
variables that occur free in the right-hand sides of the bindings of
$\valbind$, with ``free'' and ``right-hand side'' as defined in
Section~\ref{sec:appendix-derived-forms-let}; the variables bound by
$\valbind$ are the domain of the value environment it evaluates to. By the
restrictions of Section~\ref{sec:syntax-core-restrictions}, every binding of
a recursive declaration has the form
$\var\ \kw{=}\ \kw{\textbackslash}\,\atpat\ \kw{->}\ \expm$. % A26

\section{Inference Rules}\label{sec:dynamic-core-rules}

The rules infer sentences of the forms
\[
  \jdg{E}{\mathit{phrase}}{A},\qquad
  \jdgv{E}{v}{\mathit{phrase}}{A},\qquad
  \japp{v_1}{v_2}{v}
\]
where $E$ is an environment, $v$ is a value against which the
phrase is matched ($r$, a record, for pattern rows), and $A$ is the result: a
value, a record, a value environment, or FAIL (a pattern row gives a value
environment, never FAIL). The first form is pronounced
``$\mathit{phrase}$ evaluates to $A$ in $E$'', the second
``$\mathit{phrase}$, matched against $v$, gives $A$ in $E$'', and the third
``$v_1$ applied to $v_2$ gives $v$''. The hypotheses of a rule that are not
sentences are its \emph{side conditions}. The convention for
options is the same as for the Core static semantics: in each instance of a
rule, the phrases in first options $\opt{\ }$ are either all present or all
absent. % R6

Sentences have no states, and there are no exception packets: the Core has
neither assignable references nor exceptions. Hence no state convention and
no exception convention is needed, and every rule is given in full. % D2
In particular the order of the premises of a rule carries no meaning. A rule
says which results the parts of a phrase must have, and the conclusion holds
only when all of them do; since nothing but the result is observable, an
implementation may evaluate the parts of a phrase in any order. % A57, A58
Only the parts that some premise names are evaluated: the body of a lambda is
evaluated when the closure is applied, and a \kw{case} evaluates the body of
the chosen match rule only.

We allow compound variables, built from the variables of
Figures~\ref{fig:dynamic-core-simple} and~\ref{fig:dynamic-core-compound}
with ``/'', to range over unions of classes: $v/\FAIL$ ranges over
$\Val\cup\FAILcls$ and $\VE/\FAIL$ over $\ValEnv\cup\FAILcls$. We extend
modification to failure:
\[
  \VE+\FAIL \;=\; \FAIL+\VE \;=\; \FAIL+\FAIL \;=\; \FAIL .
\]

A phrase has no result in $E$ when no sentence with it on the left of
$\Rightarrow$ can be inferred. This is so when its evaluation does not
terminate, and when it needs an application of \kw{Debug.todo}. % A25, A62, audit X12
The rules place no bound on the resources evaluation uses; an implementation
that stops for lack of resources is outside the definition. % A61

\rulegroup{Atomic Expressions}{\jdg{E}{\atexp}{v}}

\defnrule[r:d-atexp-scon]{}{\jdg{E}{\scon}{\valof(\scon)}}

\defnrule[r:d-atexp-var]{E(\longvar)=v}{\jdg{E}{\longvar}{v}}

\defnrule[r:d-atexp-con]{E(\longcon)=v}{\jdg{E}{\longcon}{v}}

\defnrule[r:d-atexp-op]{}{\jdg{E}{\kw{(}\opv\kw{)}}{\OTdyn(\opv)}}

\defnrule[r:d-atexp-unit]{}{\jdg{E}{\kw{()}}{()}}

\defnrule[r:d-atexp-tuple]
  {\jdg{E}{\expm_1}{v_1}\premsep\jdg{E}{\expm_2}{v_2}\premsep
   \opt{\jdg{E}{\expm_3}{v_3}}}
  {\jdg{E}{\kw{(}\expm_1\kw{,}\ \expm_2\ \opt{\kw{,}\ \expm_3}\kw{)}}
       {(v_1,v_2\opt{,v_3})}}

\defnrule[r:d-atexp-nil]{}{\jdg{E}{\kw{[}\ \kw{]}}{\kw{[]}}}

\defnrule[r:d-atexp-record]
  {\opt{\jdg{E}{\exprow}{r}}}
  {\jdg{E}{\kw{\kwlb}\ \opt{\exprow}\ \kw{\kwrb}}{\emptymap\opt{+\,r}}}

\defnrule[r:d-atexp-update]
  {E(\var)=r\premsep\jdg{E}{\exprow}{r'}}
  {\jdg{E}{\kw{\kwlb}\ \var\ \kw{|}\ \exprow\ \kw{\kwrb}}{r+r'}}

\defnrule[r:d-atexp-paren]{\jdg{E}{\expm}{v}}{\jdg{E}{\kw{(}\expm\kw{)}}{v}}

\defnrule[r:d-atexp-let]
  {\jdg{E}{\dec}{\VE}\premsep\jdg{E+\VE}{\expm}{v}}
  {\jdg{E}{\kw{let}\ \dec\ \kw{in}\ \expm}{v}}

\begin{rulecomments}
\rcomment{r:d-atexp-scon} The value of an integer constant depends on the
type of its occurrence after number specialisation
(Section~\ref{sec:dynamic-core-numbers}). Thus, with \kw{x :\ Float} defined
by \kw{x = 2}, the constant denotes a floating-point number and
\kw{x \textasciicircum{} -1} evaluates to \kw{0.5}; and with
\kw{inc n = n + 1}, the application \kw{inc 0.5} uses the default copy of
\kw{inc}, at $\tFloat$, in which \kw{1} denotes a floating-point number, and
evaluates to \kw{1.5}, while \kw{inc (7 // 2)} uses its copy at $\tInt$ and
evaluates to the Int value $4$. % A87
\rcomment{r:d-atexp-var} As in the static semantics, long variables are
looked up in the environment. The identifier status of the static semantics
plays no part.
\rcomment{r:d-atexp-con} A constructor of a custom type is bound to the
constructor value $\con$, awaiting all its arguments; the value bound to the
constructor of a record type alias is given in
Chapter~\ref{ch:dynamic-modules}.
\rcomment{r:d-atexp-op} The value of an operator does not depend on $E$, so
a variable that a program declares cannot capture an operator: a top-level
\kw{add} of a module leaves \kw{(+)} unchanged. The operator values
\kw{(\&\&)} and \kw{(||)} denote
functions of the initial basis, which receive both arguments as values: in
\kw{(\&\&) False e} the expression \kw{e} is evaluated. Only the infixed forms
\kw{e1 \&\& e2} and \kw{e1 || e2}, which are derived forms
(Appendix~\ref{ch:appendix-derived-forms}), evaluate their second operand
conditionally. % A56
\rcomment{r:d-atexp-update} By elaboration, $\Dom r'\subseteq\Dom r$, so the
result has the fields of $r$, with the values of the fields named in
$\exprow$ replaced. The record bound to $\var$ is unaffected, there being no
state.
\rcomment{r:d-atexp-let} Every binding of $\dec$ is evaluated, whether or not
$\expm$ uses it: \kw{let \_ = Debug.todo "x" in 1} has no value. A \kw{let}
binds a single group of bindings (Appendix~\ref{ch:appendix-derived-forms}),
so the evaluation of the groups of a \kw{let} as written is ordered by the
nesting of the equivalent form, and within a group the order does not
matter. % A57
\end{rulecomments}

\rulegroup{Expression Rows}{\jdg{E}{\exprow}{r}}

\defnrule[r:d-exprow-field]
  {\jdg{E}{\expm}{v}\premsep\opt{\jdg{E}{\exprow}{r}}}
  {\jdg{E}{\lab\ \kw{=}\ \expm\ \opt{\kw{,}\ \exprow}}{\{\lab\mapsto v\}\opt{+\,r}}}

\begin{rulecomments}
\rcomment{r:d-exprow-field} The fields are evaluated in no specified order,
neither as written nor by label (Section~\ref{sec:dynamic-core-rules}). % A58
\end{rulecomments}

\rulegroup{Expressions}{\jdg{E}{\expm}{v}}

\defnrule[r:d-exp-atexp]{\jdg{E}{\atexp}{v}}{\jdg{E}{\atexp}{v}}

\defnrule[r:d-exp-app]
  {\jdg{E}{\expm}{v_1}\premsep\jdg{E}{\atexp}{v_2}\premsep\japp{v_1}{v_2}{v}}
  {\jdg{E}{\expm\ \atexp}{v}}

\defnrule[r:d-exp-lambda]
  {}
  {\jdg{E}{\kw{\textbackslash}\,\atpat\ \kw{->}\ \expm}{(\atpat,\expm,E,\emptymap)}}

\defnrule[r:d-exp-case]
  {\jdg{E}{\expm}{v}\premsep\jdgv{E}{v}{\match}{v'}}
  {\jdg{E}{\kw{case}\ \expm\ \kw{of}\ \match}{v'}}

\begin{rulecomments}
\rcomment{r:d-exp-app} The argument is evaluated whether or not the
function uses it: \kw{(\textbackslash\_ -> 1) (Debug.todo "x")} has no
value.
\rcomment{r:d-exp-lambda} The fourth component of the closure is empty
because a lambda introduces no recursively defined values;
rule~\ruleref{r:d-dec-rec} supplies it for a recursive declaration. The
closure holds the environment $E$ of this evaluation, so closures made in
successive calls of a function hold the values of their own calls. Two
evaluations of one lambda give closures that no rule distinguishes from one
another by identity. % A63, A50
\rcomment{r:d-exp-case} There is no rule in which $\match$ gives FAIL: the
static semantics requires every match of a \kw{case} to be exhaustive
(Chapter~\ref{ch:static-core}), so for a value of the right type some match
rule applies.
\end{rulecomments}

\phantomsection\label{sec:dynamic-core-apply}
\rulegroup{Application of Values}{\japp{v_1}{v_2}{v}}

\defnrule[r:d-app-closure]
  {\jdgv{E+\Rec\VE}{v}{\atpat}{\VE'}\premsep
   \jdg{E+\Rec\VE+\VE'}{\expm}{v'}}
  {\japp{(\atpat,\expm,E,\VE)}{v}{v'}}

\defnrule[r:d-app-con]
  {}
  {\japp{(\con\ v_1\cdots v_k)}{v}{\con\ v_1\cdots v_k\ v}}

\defnrule[r:d-app-basic-partial]
  {k+1<\arity(b)}
  {\japp{(b,v_1\cdots v_k)}{v}{(b,v_1\cdots v_k\,v)}}

\defnrule[r:d-app-basic-apply]
  {k+1=\arity(b)\premsep v'\in\APPLY(b,v_1\cdots v_k\,v)}
  {\japp{(b,v_1\cdots v_k)}{v}{v'}}

\begin{rulecomments}
\rcomment{r:d-app-closure} There is no rule in which $\atpat$, matched
against $v$, gives FAIL: the static semantics requires the pattern of a
lambda to be exhaustive (Chapter~\ref{ch:static-core}). Were it to fail, no
rule would apply.
\rcomment{r:d-app-con} No side condition $k<\arity(\con)$ is needed: in a
program that elaborates a constructor value is applied only while it awaits
arguments. Thus \kw{Just} applied to \kw{1} gives the complete constructor
value \kw{Just 1}.
\rcomment{r:d-app-basic-apply} The result is any member of the set
$\APPLY(b,v_1\cdots v_k\,v)$
(Section~\ref{sec:dynamic-core-basic}). When the set has two members the rule
does not determine the result: this happens only when an equality test
reaches two functions, as in \kw{(\textbackslash x -> x) == (\textbackslash
x -> x)}, which may evaluate to \kw{True} or to \kw{False}. When the set is
empty, as for \kw{Debug.todo}, the rule never applies, and the application
has no result. % A51, A62, R9
\end{rulecomments}

No rule applies the unit value, a special value, a complete constructor
value, a record, a tuple or an effect value: in a program that elaborates the
function of an application is always a closure, a constructor value awaiting
arguments, or a basic value awaiting arguments. The rules of this group are
the only definition of application; Chapters~\ref{ch:effects}
and~\ref{ch:programs} use the sentence $\japp{v_1}{v_2}{v}$ to apply function
values outside any phrase.

\rulegroup{Matches}{\jdgv{E}{v}{\match}{v'/\FAIL}}

\defnrule[r:d-match-first]
  {\jdgv{E}{v}{\mrule}{v'}}
  {\jdgv{E}{v}{\mrule\ \opt{\kw{;}\ \match}}{v'}}

\defnrule[r:d-match-fail]
  {\jdgv{E}{v}{\mrule}{\FAIL}}
  {\jdgv{E}{v}{\mrule}{\FAIL}}

\defnrule[r:d-match-next]
  {\jdgv{E}{v}{\mrule}{\FAIL}\premsep\jdgv{E}{v}{\match}{v'/\FAIL}}
  {\jdgv{E}{v}{\mrule\ \kw{;}\ \match}{v'/\FAIL}}

\begin{rulecomments}
\rcomment{r:d-match-first} The value being matched stands to the left of
the turnstile. The match rules are tried in the
order written, and the first whose pattern matches gives the result; the
match rules after it are not evaluated, even when their patterns would also
match. This order is part of the meaning, and is unaffected by the freedom of
evaluation order (Section~\ref{sec:dynamic-core-rules}), since a rule for a
later match rule requires every earlier one to give FAIL.
\end{rulecomments}

\rulegroup{Match Rules}{\jdgv{E}{v}{\mrule}{v'/\FAIL}}

\defnrule[r:d-mrule-match]
  {\jdgv{E}{v}{\pat}{\VE}\premsep\jdg{E+\VE}{\expm}{v'}}
  {\jdgv{E}{v}{\pat\ \kw{->}\ \expm}{v'}}

\defnrule[r:d-mrule-fail]
  {\jdgv{E}{v}{\pat}{\FAIL}}
  {\jdgv{E}{v}{\pat\ \kw{->}\ \expm}{\FAIL}}

\begin{rulecomments}
\rcomment{r:d-mrule-fail} The body $\expm$ is not evaluated, so the result
is FAIL even when $\expm$ has no value.
\end{rulecomments}

\rulegroup{Declarations}{\jdg{E}{\dec}{\VE}}

\defnrule[r:d-dec-nonrec]
  {\jdg{E}{\valbind}{\VE}\premsep\Dom\VE\cap\dcoreFV(\valbind)=\emptyset}
  {\jdg{E}{\valbind}{\VE}}

\defnrule[r:d-dec-rec]
  {\jdg{E}{\valbind}{\VE}\premsep\Dom\VE\cap\dcoreFV(\valbind)\neq\emptyset}
  {\jdg{E}{\valbind}{\Rec\VE}}

\begin{rulecomments}
\rcomment{r:d-dec-nonrec} Rules~\ruleref{r:d-dec-nonrec}
and~\ruleref{r:d-dec-rec} cover the one production, distinguished by whether the declaration is recursive
(Section~\ref{sec:dynamic-core-closures}). The right-hand sides are evaluated
in $E$, without the variables the declaration binds; in a recursive
declaration every right-hand side is a lambda, and $\Rec$ gives each of its
closures the bindings of the whole declaration. $\Rec$ is not applied to a
declaration that is not recursive, since a closure it binds may have been
made elsewhere, in an environment where a variable spelt as one that the
declaration binds denotes something else. A \kw{let} declaration needs only
$\Rec$ because every binding of a recursive \kw{let} declaration is a
function (Section~\ref{sec:syntax-core-restrictions}); top-level declarations,
whose recursive members need not be functions, are treated in
Chapter~\ref{ch:dynamic-modules}. % A26
\end{rulecomments}

\rulegroup{Value Bindings}{\jdg{E}{\valbind}{\VE}}

\defnrule[r:d-valbind-var]
  {\jdg{E}{\expm}{v}\premsep\opt{\jdg{E}{\valbind}{\VE}}}
  {\jdg{E}{\var\ \kw{=}\ \expm\ \opt{\kw{and}\ \valbind}}{\{\var\mapsto v\}\opt{+\,\VE}}}

\defnrule[r:d-valbind-pat]
  {\jdg{E}{\expm}{v}\premsep\jdgv{E}{v}{\atpat}{\VE}}
  {\jdg{E}{\atpat\ \kw{=}\ \expm}{\VE}}

\begin{rulecomments}
\rcomment{r:d-valbind-pat} There is no rule in which $\atpat$, matched
against $v$, gives FAIL: the static semantics requires the pattern of a value
binding to be exhaustive (Chapter~\ref{ch:static-core}).
\end{rulecomments}

\rulegroup{Atomic Patterns}{\jdgv{E}{v}{\atpat}{\VE/\FAIL}}

\defnrule[r:d-atpat-wild]{}{\jdgv{E}{v}{\kw{\_}}{\emptymap}}

\defnrule[r:d-atpat-scon]
  {v=\valof(\scon)}
  {\jdgv{E}{v}{\scon}{\emptymap}}

\defnrule[r:d-atpat-scon-fail]
  {v\neq\valof(\scon)}
  {\jdgv{E}{v}{\scon}{\FAIL}}

\defnrule[r:d-atpat-var]{}{\jdgv{E}{v}{\var}{\{\var\mapsto v\}}}

\defnrule[r:d-atpat-con]
  {E(\longcon)=\con\premsep v=\con}
  {\jdgv{E}{v}{\longcon}{\emptymap}}

\defnrule[r:d-atpat-con-fail]
  {E(\longcon)=\con\premsep v\neq\con}
  {\jdgv{E}{v}{\longcon}{\FAIL}}

\defnrule[r:d-atpat-unit]{}{\jdgv{E}{()}{\kw{()}}{\emptymap}}

\defnrule[r:d-atpat-tuple]
  {v=(v_1,v_2\opt{,v_3})\premsep
   \jdgv{E}{v_1}{\pat_1}{\VE_1/\FAIL}\\
   \jdgv{E}{v_2}{\pat_2}{\VE_2/\FAIL}\premsep
   \opt{\jdgv{E}{v_3}{\pat_3}{\VE_3/\FAIL}}}
  {\jdgv{E}{v}{\kw{(}\pat_1\kw{,}\ \pat_2\ \opt{\kw{,}\ \pat_3}\kw{)}}
        {\VE_1/\FAIL+\VE_2/\FAIL\opt{+\,\VE_3/\FAIL}}}

\defnrule[r:d-atpat-nil]{v=\kw{[]}}{\jdgv{E}{v}{\kw{[}\ \kw{]}}{\emptymap}}

\defnrule[r:d-atpat-nil-fail]{v\neq\kw{[]}}{\jdgv{E}{v}{\kw{[}\ \kw{]}}{\FAIL}}

\defnrule[r:d-atpat-record]
  {v=r\premsep\opt{\jdgv{E}{r}{\patrow}{\VE}}}
  {\jdgv{E}{v}{\kw{\kwlb}\ \opt{\patrow}\ \kw{\kwrb}}{\emptymap\opt{+\,\VE}}}

\defnrule[r:d-atpat-paren]
  {\jdgv{E}{v}{\pat}{\VE/\FAIL}}
  {\jdgv{E}{v}{\kw{(}\pat\kw{)}}{\VE/\FAIL}}

\begin{rulecomments}
\rcomment{r:d-atpat-scon} A constant in a pattern is an
integer, character or string constant, never a floating-point one, so the
test $v=\valof(\scon)$ is identity of special values. An integer constant
matches the normalised integer it denotes. % A7, A9
\rcomment{r:d-atpat-var} A variable pattern always binds: constructors and
variables are spelt differently (Chapter~\ref{ch:syntax-core}), so no
identifier status is consulted.
\rcomment{r:d-atpat-unit} For rules~\ruleref{r:d-atpat-unit},
\ruleref{r:d-atpat-tuple} and~\ruleref{r:d-atpat-record}: in a program that elaborates the value matched has the
form the pattern requires, so no rule for these patterns gives FAIL on its
own account; a tuple fails only when a component pattern fails. A record
pattern matches a record with any further fields.
\end{rulecomments}

\rulegroup{Pattern Rows}{\jdgv{E}{r}{\patrow}{\VE}}

\defnrule[r:d-patrow-field]
  {\lab\in\Dom r\premsep\opt{\jdgv{E}{r}{\patrow}{\VE}}}
  {\jdgv{E}{r}{\lab\ \opt{\kw{,}\ \patrow}}{\{\lab\mapsto r(\lab)\}\opt{+\,\VE}}}

\begin{rulecomments}
\rcomment{r:d-patrow-field} In the result, $\lab$ stands for the variable
spelt as the label, which the pattern row binds. The side condition
$\lab\in\Dom r$ makes the rule apply exactly when $r$ has the field; by
elaboration it always holds, since a record of the type the pattern requires
has every field the pattern row names. Hence a pattern row cannot fail, and
its result is a value environment, never FAIL, as
rule~\ruleref{r:d-atpat-record} expects. % D5, mechanisation F3, F5
\end{rulecomments}

\rulegroup{Patterns}{\jdgv{E}{v}{\pat}{\VE/\FAIL}}

\defnrule[r:d-pat-atpat]
  {\jdgv{E}{v}{\atpat}{\VE/\FAIL}}
  {\jdgv{E}{v}{\atpat}{\VE/\FAIL}}

\defnrule[r:d-pat-con]
  {E(\longcon)=\con\premsep v=\con\ v_1\cdots v_n\\
   \jdgv{E}{v_1}{\atpat_1}{\VE_1/\FAIL}\ \ \cdots\ \
   \jdgv{E}{v_n}{\atpat_n}{\VE_n/\FAIL}}
  {\jdgv{E}{v}{\longcon\ \atpat_1\cdots\atpat_n}
        {\VE_1/\FAIL+\cdots+\VE_n/\FAIL}}

\defnrule[r:d-pat-con-fail]
  {E(\longcon)=\con\premsep v\notin\{\con\}\times\Val^n}
  {\jdgv{E}{v}{\longcon\ \atpat_1\cdots\atpat_n}{\FAIL}}

\defnrule[r:d-pat-cons]
  {v=v_1\mathbin{\kw{::}}v_2\\
   \jdgv{E}{v_1}{\pat_1}{\VE_1/\FAIL}\premsep
   \jdgv{E}{v_2}{\pat_2}{\VE_2/\FAIL}}
  {\jdgv{E}{v}{\pat_1\ \kw{::}\ \pat_2}{\VE_1/\FAIL+\VE_2/\FAIL}}

\defnrule[r:d-pat-cons-fail]
  {v=\kw{[]}}
  {\jdgv{E}{v}{\pat_1\ \kw{::}\ \pat_2}{\FAIL}}

\defnrule[r:d-pat-as]
  {\jdgv{E}{v}{\pat}{\VE/\FAIL}}
  {\jdgv{E}{v}{\pat\ \kw{as}\ \var}{\VE/\FAIL+\{\var\mapsto v\}}}

\begin{rulecomments}
\rcomment{r:d-pat-con} Here $n\geq 1$, and by elaboration $n=\arity(\con)$.
All the argument patterns are matched, and the result is FAIL when any of
them fails. Since matching has no effect and always has a result, the order
in which the arguments are matched is immaterial.
\rcomment{r:d-pat-con-fail} Rules~\ruleref{r:d-atpat-scon-fail},
\ruleref{r:d-atpat-con-fail}, \ruleref{r:d-atpat-nil-fail},
\ruleref{r:d-pat-con-fail} and~\ruleref{r:d-pat-cons-fail} are the only rules
that give FAIL without a premise giving FAIL; every other FAIL comes from
one of them.
\rcomment{r:d-pat-as} By the syntactic restrictions, $\var$ is not bound by
$\pat$, so the modification adds a binding and replaces none.
\end{rulecomments}
